\begin{document}

\loomcover
  {信息论}
  {第二讲 · 从信源符号到可译字符串\\[2pt]{\small\color{selvage}Oxford B8.4 · 2024--25}}
  {依据手写笔记整理的读填笔记}
  {\today}

\section*{一条主线与知识地图}
\addcontentsline{toc}{section}{一条主线与知识地图}
\warmth{0}

\begin{strand}
信源编码把每个信源符号替换成一个有限字符串。编码质量有两个侧面：以期望码长衡量的
\keyword{效率}，以及在码字拼接后检验的\keyword{可译性}。核心设计问题是在缩短字符串的
同时，保留足以识别码字边界的结构。
\end{strand}

\block{目前的知识地图}
\noindent\begin{tabularx}{\linewidth}{@{}c l L@{}}
\toprule
{\color{indigo}站} & {\color{indigo}对象} & \multicolumn{1}{l}{\color{madder}要回答的问题}\\
\midrule
1 & $C:\X\to\D^*$ & 每个信源符号由哪个有限字符串表示？\\
2 & $L(C)=\E[\ell(X)]$ & 平均要使用多少个编码符号？\\
3 & $C^*:\X^*\to\D^*$ & 码字拼接起来以后还能否唯一解码？\\
\bottomrule
\end{tabularx}

\block{全文约定}
信源字母表 $\X$ 有限，$\D^*$ 表示由 $D$ 元编码字母表 $\D$ 中的符号组成的所有
有限长字符串。例子通常取 $\D=\{0,1\}$。

\clearpage
{\small\tableofcontents}

\clearpage
\section*{第二讲 · 信源编码}
\addcontentsline{toc}{section}{第二讲 · 信源编码}
\begin{strand}
先给信源符号分配字符串，再按照信源分布计算平均长度，最后提出决定性问题：拼接操作
是否保留了足够的结构，使我们能够逆转编码过程？
\end{strand}

% ===========================================================================
\section{编码、码长与可译性}\label{sec:codes}
\warmth{0}\quad\whisper{一个码字标记一个符号；扩展码还必须保住整个符号序列。}

\begin{strand}
编码从一个“信源符号到有限字符串”的映射开始。短字符串能够节省空间，但不同的符号，
乃至不同的信源序列，都必须保持可区分。因此，期望码长与可译性必须放在一起研究。
\end{strand}

\block{从信源符号到码字}

\begin{definition}[信源编码]\label{def:source-code}
设随机变量 $X$ 的有限字母表为 $\X$。一个 $D$ 元信源编码是映射
\[
  C:\X\longrightarrow\D^*,
\]
其中 $\D^*$ 是 $D$ 元字母表 $\D$ 上所有\answerin[3.5cm]{有限长字符串}的集合。
字符串 $C(x)$ 称为 $x$ 的\keyword{码字}，$|C(x)|$ 是它的码长。
\studentrecall{空中应填“有限长字符串”。}
\end{definition}

例如，若 $\X=\{\text{红},\text{蓝}\}$ 且 $\D=\{0,1\}$，可以取
$C(\text{红})=00$、$C(\text{蓝})=11$。

\begin{definition}[期望码长]\label{def:expected-length}
若 $X$ 的概率质量函数为 $p$，编码 $C$ 的期望码长定义为
\[
  L(C)=\E[|C(X)|]
      =\sum_{x\in\X}p(x)|C(x)|.
\]
\end{definition}

\begin{example}[与熵匹配的二进制编码]\label{ex:dyadic-code}
令 $\X=\{1,2,3,4\}$，其概率和码字如下：
\[
\begin{array}{c|cccc}
x&1&2&3&4\\ \hline
p(x)&\tfrac12&\tfrac14&\tfrac18&\tfrac18\\
C(x)&0&10&110&111
\end{array}
\]
四个码长分别为 $1,2,3,3$，所以
\[
  L(C)=\tfrac12\cdot1+\tfrac14\cdot2
       +\tfrac18\cdot3+\tfrac18\cdot3
       =\answerin[1.4cm]{\tfrac74}\ \text{比特}.
\]
这里每个符号都满足 $|C(x)|=-\log_2p(x)$，因此
\[
  \Ent(X)=\sum_xp(x)\log_2\frac1{p(x)}
         =\answerin[1.4cm]{\tfrac74}\ \text{比特}=L(C).
\]
\studentrecall{两个空都应填 $\tfrac74$ 比特。}
\end{example}

\begin{yourturn}
为什么这个例子中的 $L(C)=\Ent(X)$ 不需要特殊的平均技巧？
\studentrecall{每个 $x$ 都满足 $|C(x)|=-\log_2p(x)$，所以等式逐项成立。}
\ifshowanswers\else\workspace[2]\fi
\answerprose{这个等式对每一项都成立：每个码长都满足
$|C(x)|=-\log_2p(x)$。两边乘以 $p(x)$ 后求和，立即得到
$L(C)=\Ent(X)$。}
\end{yourturn}

\block{三个层次的可译性}

\begin{definition}[非奇异码]\label{def:nonsingular}
若不同信源符号具有不同码字，则称编码 $C$ 为\keyword{非奇异码}：
\[
  x\ne x'\quad\Longrightarrow\quad C(x)\ne C(x').
\]
等价地说，$C$ 是一个\answerin[2.2cm]{单射}。
\studentrecall{空中应填“单射”。}
\end{definition}

非奇异性只保护单个符号。为了编码符号序列，将 $C$ 按照拼接方式扩展为
\[
  C^*:\X^*\longrightarrow\D^*,\qquad
  C^*(x_1x_2\cdots x_n)=C(x_1)C(x_2)\cdots C(x_n).
\]

\begin{definition}[唯一可译码]\label{def:ud}
若扩展码 $C^*$ 是非奇异的，则称 $C$ 为\keyword{唯一可译码}。换言之，每个编码后的
字符串至多对应一个产生它的信源序列。
\end{definition}

\begin{definition}[前缀码]\label{def:prefix}
若任何码字都不是另一不同码字的前缀，则称该编码为\keyword{前缀码}，也称即时码。
\end{definition}

每个前缀码都是唯一可译码，每个唯一可译码都是非奇异码：
\[
  \{\text{前缀码}\}
  \subset
  \{\text{唯一可译码}\}
  \subset
  \{\text{非奇异码}\}.
\]
反方向一般不成立。

\begin{example}[唯一可译但不是前缀码]\label{ex:ud-not-prefix}
对 $\X=\{A,B,C,D\}$，考虑
\[
  C(A)=0,\qquad C(B)=01,\qquad C(C)=011,\qquad C(D)=111.
\]
它不是前缀码，因为 $0$ 是 $01$ 和 $011$ 的前缀；但它仍是唯一可译码。从较长码字
中删去开头的 $0$，只会留下 $1$ 或 $11$；再将这些后缀与 $111$ 比较，只会使
$1\leftrightarrow11$ 来回转换，永远不会得到空后缀。因此不会有两个信源序列产生
相同的拼接结果。
\end{example}

\begin{yourturn}
编码 $\{0,01,10\}$ 是非奇异码，却不是唯一可译码。请给出两个编码结果相同的不同
信源序列。
\studentrecall{$AC$ 与 $BA$ 都编码为 $010$。}
\ifshowanswers\else\workspace[1]\fi
\answerprose{$AC$ 与 $BA$ 都编码为 $010$。}
\end{yourturn}

\vspace{1.2cm}
\recall{为什么非奇异性不能保证拼接之后仍可译？}


\end{document}
